
\paragraph{Algebraic steppers.}
Algebraic steppers for functional languages have long served pedagogy. The
DrScheme/Racket algebraic stepper presents small-step reductions that mirror
operational semantics, with formal correctness arguments for its
traces~\cite{clements2001stepper}. A parallel rewriting-based tradition runs
through the HIPE/WinHIPE line, in which user programs are evaluated by an
exposed rewriting system and the resulting traces are animated for classroom
use~\cite{urquiza2007winhipe}. Cong and Asai showed how delimited continuations
can implement steppers for ML-like languages~\cite{cong2016stepper}, a line
continued in steppers for practical OCaml fragments that handle effectful
constructs such as exceptions and printing
primitives~\cite{furukawa2019stepping}, and most recently extended to simple
modules with hierarchical variable references~\cite{asai2025modules}. Closest
to our motivation, Olmer, Heeren, and Jeuring's Haskell tutoring environment
addresses trace verbosity by letting students switch between innermost and
outermost evaluation orders to control which reductions are
shown~\cite{olmer2014evaluating}; the surrounding Ask-Elle tutor integrates
this with automated feedback on student-written Haskell
programs~\cite{gerdes2017askelle}. Our work keeps the source-level,
semantics-directed view of evaluation common to all of these systems, but
gives users a finer-grained control: rather than picking a strategy globally,
users embed lexically scoped patterns that select which reductions appear in
the trace on a redex-by-redex basis.

\paragraph{Programmer-directed focus.}
Bertot's occurrence-based debugger specifications and Bernstein and Stark's
focusing debugger formalize how users can direct debugging attention to
particular source-level
occurrences~\cite{bertot1991occurrences,bernstein1995focusing}. Algorithmic
debugging, going back to Shapiro~\cite{shapiro1983algorithmic} and continuing
through modern Haskell instantiations such as Buddha and
Hoed~\cite{pope2004buddha,faddegon2016hoed}, achieves user-directed focus
through oracle questioning: the programmer's yes/no answers about
subcomputation correctness successively narrow the debugger's attention to
the buggy subtree. Closely related, Perera et al.\ recast a functional
evaluation trace as an \emph{explanation} of an observed output, using
program slicing to extract the minimal subprogram that produced
it~\cite{perera2012explain}. Interrogative debugging via the
Whyline takes a complementary approach, generating a menu of ``why did'' and
``why didn't'' questions that the user selects to highlight relevant runtime
events~\cite{ko2004whyline}. Expositor exposes execution traces as first-class
values that programmers can script with \texttt{map} and \texttt{filter} to
direct focus~\cite{khoo2013expositor}. On the imperative side, Boothe's
bidirectional debugger gave programmers an event-counter-based language for
navigating execution forwards and backwards along chosen event
classes~\cite{boothe2000bidirectional}. Hazel filters extend this lineage with
redex-structural patterns that compose through lexical scope and integrate
with an algebraic stepper.

\paragraph{Lexically scoped source annotations.}
The methodological pattern of attaching lightweight, lexically scoped
annotations to source code to control a dynamic observation has a direct
ancestor in profiling. Sansom and Peyton Jones introduced \emph{cost centres}
for higher-order functional languages: programmer-written annotations,
attached to source regions, that attribute runtime costs to named scopes,
with a formal cost-attribution semantics that is invariant under
evaluation~\cite{sansom1997profiling}. Filter annotations share this
methodological shape. A lexically scoped construct in the source does not
affect the result of evaluation, but governs how a dynamic event (cost
attribution there, step visibility here) is attributed to source regions.

\paragraph{Tracing tools for functional languages.}
Tracing tools record the reduction history of a computation for later
inspection. An early source-level debugger from Standard ML by Tolmach and
Appel used source-to-source instrumentation and first-class continuations to
provide breakpoints and reverse execution without depending on the machine
representation~\cite{tolmach1995debugger}. Nilsson and Fritzson's Freja
debugger introduced declarative debugging for lazy functional languages, using
an evaluation-dependence tree as the basis for oracle
questioning~\cite{nilsson1994algorithmic}. Hood lets programmers insert
\texttt{observe} annotations to capture intermediate values without altering
program semantics~\cite{gill2000hood}, and GHood adds graphical animation of
those observations~\cite{reinke2001ghood}. Hat records an \emph{augmented
redex trail} and provides multiple views over it, including search and
structured querying~\cite{wallace2001hat}. Chitil's source-based trace
explorer builds on Hat to unify algorithmic debugging, stepping, and dynamic
slicing into a single source-driven viewer in which the user navigates the
trace by selecting source-level subterms~\cite{chitil2005source}. The GHCi
debugger brought interactive source-level breakpoints to lazy Haskell via
lightweight bytecode instrumentation~\cite{marlow2007ghci}. More recently,
CSI:~Haskell extends GHC's coverage tooling to record evaluation traces
alongside the call stack, reconstructing the chain of events leading to a
fault in lazy code~\cite{gissurarson2023csi}. Our stepper similarly maintains
a step history, but the trace is \emph{filtered} by user-specified patterns
rather than recorded uniformly.

\paragraph{Pattern languages and term-rewriting strategies.}
Our filter language is also related to program-pattern and
strategy-controlled rewriting systems. SCRUPLE and ASTLOG use structural
patterns to query program trees~\cite{paul1994scruple,crew1997astlog}.
Stratego separates rewrite rules from strategies that control where and how
rules are applied~\cite{visser2001stratego}. Coccinelle's \emph{semantic
patches} generalize Unix-style patches to a structural pattern language over
C, and have driven thousands of automated evolutions in the Linux
kernel~\cite{padioleau2008coccinelle}. Hazel filters use a smaller pattern
language and do not let users change the evaluator. Instead, patterns control
the visibility of steps produced by a fixed algebraic stepper.

\paragraph{Programming education context.}
More broadly, our stepper sits within a long-running line of work on program
visualization and \emph{notional machines} for introductory programming
education. Sorva, Karavirta, and Malmi survey generic visualization systems
and distill the design lessons of that body of work, including the central
tension between fidelity to the underlying semantics and clarity for a
novice~\cite{sorva2013review}. Clements and Krishnamurthi's recent
investigation of pedagogical models for runtime stacks and scope --- using a
block-based programming environment for empirical study --- illustrates that
the question of how best to render evaluation comprehensible to students
remains an active research concern~\cite{clements2022notional}.

\paragraph{Hazel substrate.}
Hazelnut-style structural editing and semantics for incomplete programs make
the AST explicit at all times, which naturally supports tree-pattern matching
and stepping in the presence of holes~\cite{omar2017hazelnut,omar2019live}.
The substrate has since been extended with typed pattern holes that enable
live evaluation in the presence of incomplete match
expressions~\cite{yuan2023peanut}, and with the marked lambda calculus that
gives a total account of type error localization and recovery, providing the
foundation on which modern Hazel is built~\cite{zhao2024marking}. The filter
calculus builds on this substrate to give a precise semantics for scoped,
pattern-based stepping control in a live functional environment.
